\chapter{Portal Vein Thrombosis}
\index{Portal Vein Thrombosis}

\medskip
\noindent\textit{This chapter is for general education. A suspected blood clot needs prompt medical assessment.}

\section*{What it is}
Portal vein thrombosis is a blood clot that partly or completely blocks the portal vein, which carries blood from the digestive organs to the liver. It may be recent (acute) or longstanding (chronic). Causes include cirrhosis, abdominal infection or inflammation, cancer, surgery or injury, pregnancy, and conditions that make blood clot more easily; sometimes no cause is found.

\section*{Symptoms}
Some people have no symptoms. Possible features include abdominal pain, fever, nausea, an enlarged spleen, or fluid in the abdomen. Chronic blockage can cause portal hypertension, enlarged veins in the oesophagus or stomach, and gastrointestinal bleeding. Vomiting blood, passing black stools, fainting, or sudden severe abdominal pain requires emergency care.

\section*{When to Seek Medical Care}
Seek urgent care for vomiting blood, black stools, fainting, sudden severe abdominal pain, severe weakness, or rapidly worsening symptoms.

\section*{Diagnosis}
Assessment usually includes a history, examination, blood tests, and Doppler ultrasound. CT or MRI may show the extent of the clot, complications, and an underlying cause. Clinicians may investigate liver disease, infection, cancer, and inherited or acquired clotting disorders.

\section*{Treatment}
Treatment is individualised. It may include anticoagulation, treatment of infection or inflammation, and management of cirrhosis or another cause. Bleeding from enlarged veins may require urgent endoscopic treatment and measures to lower portal pressure. Do not start aspirin, an anticoagulant, or a herbal product for this condition without medical advice, because bleeding risk and clot risk must be balanced.

\section*{Outlook and follow-up}
The outlook depends on whether the clot is acute or chronic, the liver's condition, the cause, and complications. Follow-up imaging and blood tests may be needed. Take prescribed medicines consistently and ask before using medicines that increase bleeding risk.

\section*{Sources and further reading}
\begin{itemize}
\item MSD Manual Professional Edition. \textit{Portal vein thrombosis}. \url{https://www.msdmanuals.com/professional/hepatic-and-biliary-disorders/vascular-disorders-of-the-liver/portal-vein-thrombosis}
\item National Institute of Diabetes and Digestive and Kidney Diseases. \textit{Portal hypertension}. \url{https://www.niddk.nih.gov/health-information/liver-disease/portal-hypertension}
\end{itemize}
